\begin{trapbox}
\begin{enumerate}
  \item \textbf{“任意一点”漏掉端点限制。}
  必须有 $P\ne A_1,A_2$，即 $y_0\ne0$。
  当 $P=A_1$ 或 $A_2$，第三条切线就是 $x=-a$ 或 $x=a$，
  与另一条端点切线平行；交点 $P_1,P_2$ 不能按题意同时唯一确定。
  因此不能直接代入 $y_0=0$。

  \item \textbf{把“端点切线”认成“准线”。}
  本题固定线是 $x=\pm a$；椭圆的准线是
  $x=\pm a^2/c$（$c^2=a^2-b^2$），通常不重合。
  在准线上截出的两段高度，乘积一般不等于 $b^2$。

  \item \textbf{有符号纵坐标与长度混用。}
  实际竖段长是 $d_1=|y_L|$、$d_2=|y_R|$。
  若 $P$ 在下半椭圆，$y_L,y_R<0$，不能把
  $d_i$ 写成对应负的纵坐标；正确的是
  $d_1d_2=|y_Ly_R|=b^2$。

  \item \textbf{没证明同侧就直接写面积。}
  首先由 $y_Ly_R=b^2>0$ 证明交点同侧，
  才能写 $S=a(d_1+d_2)=a|y_L+y_R|$。
  对一般两条竖线的交点，若纵坐标异号，
  原来的直角梯形模型不一定成立。

  \item \textbf{面积中的底、高看反。}
  平行边是两条\textbf{竖直}线段 $A_1P_1,A_2P_2$；
  两平行边间的距离才是 $2a$。
  正确公式是
  $S=\tfrac12(d_1+d_2)\cdot2a=a(d_1+d_2)$，
  不是 $\tfrac12d_1d_2$，也不是 $2a(d_1+d_2)$。

  \item \textbf{基本不等式忘写等号条件。}
  $d_1+d_2\ge2\sqrt{d_1d_2}=2b$，
  取等当且仅当 $d_1=d_2=b$；
  此时切线必须水平，$P=(0,\pm b)$。
  不能只写“对称时取等”而不验证动点确实存在。

  \item \textbf{认为面积有最大值。}
  当 $P$ 接近长轴端点（但不取到），
  一条竖段趋于无穷大，另一条趋于零，
  $S=a(d_1+d_2)\to+\infty$。有最小值，\textbf{没有最大值}。

  \item \textbf{切线斜率法漏掉竖直切线的处理。}
  $y=kx+m$ 不包含竖直直线；本题恰好排除了
  仅有的两个竖直切点 $A_1,A_2$，所以用此形式没有遗漏。
  如果题目改变动点范围，应重新核对。

  \item \textbf{把椭圆切线当作圆的切线直接套等长。}
  “从同一点向圆作两切线，切线段相等”是圆的几何定理；
  对一般椭圆并不成立。正确的几何证明必须先把椭圆
  横向还原为圆，再在圆中用切线长定理，最后利用
  水平伸缩保持竖直长度的性质返回椭圆。

  \item \textbf{混淆其他圆锥曲线的有符号乘积。}
  对双曲线 $x^2/a^2-y^2/b^2=1$，
  非竖直切线 $y=kx+m$ 满足 $m^2-a^2k^2=-b^2$；
  它在 $x=\pm a$ 处的\textbf{有符号}纵坐标积为 $-b^2$，
  与椭圆相反。不能不加条件地搬运梯形面积最值结论。
\end{enumerate}
\end{trapbox}
